\section{Rooks on the residual mask}\label{sec:torres}
\label{sec:rook} The reduction to an auxiliary permanent, used for the constant term~\cite{tc},
lets us compare two arrival boards. On the first board we keep the available cells and lift the
diagonal restrictions between the pieces that remain to be placed. Its completions are the perfect
matchings of the mask $M$, and their weighted sum is $\operatorname{per}M$. On the second board we
fix two rooks in cells $e,f$ on a diagonal. The remaining pieces occupy the mask $M_{ef}$, which
loses two rows and two columns. The ratio $w_ew_f\operatorname{per}M_{ef}/\operatorname{per}M$
gives the probability of occupying the pair. When we approximate both permanents, how much of the
error on the first board cancels on the second? In Figure~\ref{fig:mascara-ilustrada} one can read
off the numerator for $e=(2,4)$ and $f=(3,3)$ by hand. Only the assignment $(4,1)$ remains. Thus
one of the two matchings of the mask contains this pair.

\paragraph{The class and two operations.}
\[
  w_e\in\{0\}\cup[w_0,w_1],\quad P=\operatorname{diag}(u)M\operatorname{diag}(v),
  \quad P_e\le M_0/k,\quad \|P-\mathbb U_k\|_{\rm op}\le1-\gamma .
\]
Here $0<w_0\le w_1$, $M_0\ge1$, $0<\gamma\le1$, $M$ has total support, $P$ is its doubly
stochastic scaling, and $\mathbb U_k$ has entries $1/k$. This defines the class $\mathcal K$. In
this section $k$ is the integer size of the mask, distinct from the function $\numerr(x)$ in the
integral. Occupying $e$ removes its row and column and leaves the minor $M_e$. Forbidding $e$
retains both fibers and sets $w_e=0$ in $M\setminus e$. The two operations perturb the sums to be
rescaled in different ways. The queen child $M_e^{\rm fis}$ also forbids the cells on both
diagonals through $e$. Lemma~\ref{lem:hijo-fisico} controls the scaling during this operation.

The stationary class (H) requires a rook point with unit margins on the branch $0<p_e<1/2$,
satisfying $p_e(1-p_e)=w_e\exp(z_i+y_j)$, $p_e\le C_p/k$ and $\lambda_{\min}(S_T+\Pi)\ge\gamma_A$.
Here $S_T$ is the incidence Hessian and $\Pi$ projects onto the row-minus-column mode, as in
Appendix~\ref{app:rook}. In $\mathcal K$, from the threshold $k_{\mathcal K}(M_0,\gamma)$ onward,
this point exists and is unique, with $C_p=3M_0/2$, $\gamma_A=0.88\gamma$ and
$p_e/P_e=1+O(k^{-1})$ on the support. We compute $L_T$ with $p$ and its first correction with $P$.

\paragraph{The approximation to be transported [D].}
\[
   g_T(M)=\log\operatorname{per}M-L_T(M)
       =\frac{\coefT(kP)}k+R(M),\qquad |R(M)|\le Ck^{-2}.
\]
This estimate for the value is (T). Writing $\coefT=\mathfrak F-1/12$ and evaluating $\mathfrak F$
at the rook profile along the trajectory gives the $F(x)$ in~\eqref{eq:main}. The static estimate
$(L^*)$ applies to the doubly stochastic matrix $P$ with the two scaling parameters $M_0,\gamma$.
It gives $\mathcal E(P)=\mathfrak f_{\rm e}(kP)/k+O(k^{-2})$. Adding the expansion of
$g_T-\mathcal E(P)$ gives (T) for $k$ above both thresholds. On the common class of the static and
deterministic increment estimates we have $(T^+)$,
\[
 |R(M_e)-R(M)|\le Ck^{-3},\qquad |R(M\setminus e)-R(M)|\le Ck^{-4}.
\]
Appendix~\ref{app:rook} derives both bounds in $\mathcal K$ with constants depending on
$M_0,\gamma$. As a weight is reduced, the continuation is uniform down to zero. When applying a
further expansion on the minor, we use the threshold for its own diffusion and gap parameters.

To use (S), let $r_c(N)=L_T(N_c)-L_T(N)-\log(\pi_c^{\rm der}(N)/w_c)$ be the insertion remainder
at $c$, with $\pi_c^{\rm der}=\partial_{\log w_c}L_T$. If $N$ satisfies (H), $e,f$ occupy distinct
rows and columns, and $c$ retains both fibers in $N_{ef}$, then $|r_c(N_{ef})-r_c(N)|\le C_S/k^3$.
We apply the bound for a distant minor twice, with the thresholds for $(C_p,\gamma_A)$ at $k$ and
for $(5C_p/4,3\gamma_A/4)$ at $k-1$. This comparison of minors enters the marginal in
Appendix~\ref{app:pair}.

\paragraph{From two rooks to a pair of queens.}
On a binary mask, the exact probability of occupying $E=\{e,f\}$, with distinct rows and columns,
is $\pi^{\rm ex}=\operatorname{per}M_{ef}/\operatorname{per}M$. For a single two-cell line, the
corresponding term in the logarithm of the count is $\log(1-\pi^{\rm ex})$. The local
approximation $\pi_1$ is defined in Appendix~\ref{app:pair}. If a third cell $c$ in a different
row and column is revealed, the change in the error $\pi^{\rm ex}-\pi_1$ is at most
$C\varpi(E)/k^5$, where $\varpi(E)=k^2p_ep_f$. This statement (R1) applies to the rook minor that
retains the pair. In the physical geometry a move may also forbid additional cells. In the arrival
average, the change in the rook error contributes $P_T$. The effect of one line contributes to
$J$, and the two diagonals of a complete square also produce the term $-S$ in
Appendix~\ref{app:boundary}.
